% -----------------------------------------------------------
\section{Condescension}
% -----------------------------------------------------------
	\label{CONDESCENSION}
	
% % ----------------------
% \subsection{See also}
% % ----------------------

% ----------------------
\subsection{1828 American Dictionary of the English Language} % *1828
% ----------------------
	\label{CONDESCENSION-1828}

	% -----
	\subsubsection{Condescension}
	% -----

		\begin{compactitem}
			\item Voluntary descent from rank, dignity or just claims; relinquishment of strict right; submission to inferiors in granting requests or performing acts which strict justice does not require. Hence, courtesy.
		\end{compactitem}
		
	% -----
	\subsubsection{Condescend}
	% -----
	
		\begin{compactitem}
			\item[1] To descend from the privileges of superior rank or dignity, to do some act to an inferior, which strict justice or the ordinary rules of civility do not require. Hence, to submit or yield, as to an inferior, implying an occasional relinquishment of distinction.
			\item[2] To recede from ones rights in negotiation, or common intercourse, to do some act, which strict justice does not require.
			\item[3] To stoop or descend; to yield; to submit; implying a relinquishment of rank, or dignity of character, and sometimes a sinking into debasement.
		\end{compactitem}

% % ----------------------
% \subsection{Oxford English Dictionary} % *OED
% % ----------------------
	% \label{CONDESCENSION-OED}

	% \begin{compactitem}
		% \item[]
	% \end{compactitem}

% ----------------------
\subsection{Scriptural Usage} % *SCRIPTURES
% ----------------------
	\label{CONDESCENSION-SCRIPTURES}

	% % -----
	% \subsubsection{Old Testament} % *OT
	% % -----
		% \label{CONDESCENSION-OT}
	
		% % --
		% \paragraph{KJV}
		% % --
			% \label{CONDESCENSION-OT-KJV}
	
			% \begin{tabular}{ll}
				% Total occurrences in the OT & 
			% \end{tabular}
			
			% \begin{compactdesc}
				% \item[]
			% \end{compactdesc}
			
		% % --
		% \paragraph{Hebrew Bible}
		% % --
			% \label{CONDESCENSION-HEBREW}
		
			% %
			% \subparagraph{\texorpdfstring{\he{sample word}}{}}
			% %
				% \label{PAGA} % whatever is in step bible without periods
			
				% \begin{compactdesc}
					% \item[HALOT , PDF ] \textbf{}
						% \begin{compactitem}
							% \item[1]
						% \end{compactitem}
					% \item[BDB , PDF ] \textbf{}
						% \begin{compactitem}
							% \item[1]
						% \end{compactitem}
					% \item[DCH PDF ] \textbf{}
						% \begin{compactitem}
							% \item[1]
						% \end{compactitem}
				% \end{compactdesc}
				
				% \begin{compactdesc}
					% \item[Some OT uses] \textbf{}
						% \begin{compactdesc}
							% \item[]
						% \end{compactdesc}
				% \end{compactdesc}
			
		% % --
		% \paragraph{Septuagint}
		% % --
			% \label{CONDESCENSION-LXX}
		
			% %
			% \subparagraph{\texorpdfstring{\gr{sample word}}{}}
			% %
		
	% -----
	\subsubsection{New Testament} % *NT
	% -----
		\label{CONDESCENSION-NT}
	
		% --
		\paragraph{KJV}
		% --
			\label{CONDESCENSION-NT-KJV}
			
	
			\begin{tabular}{ll}
				Total occurrences in the NT & 1
			\end{tabular}
			
			\begin{compactdesc}
				\item[Romans 12:16] Be of the same mind one toward another.  Mind not high things, but condescend to men of low estate.  Be not wise in your own conceits.
			\end{compactdesc}
			
		% --
		\paragraph{Greek New Testament} % *GREEK
		% --
			\label{CONDESCENSION-GREEK}
		
			%
			\subparagraph{\texorpdfstring{\gr{συναπάγω}}{}}
			%
				\label{SUNAPAGO}
			
				\begin{compactdesc}
					\item[BDAG 858a] \textbf{}
						\begin{compactitem}
							\item in our literature only pasive and only figurative
							\item[1] to cause someone in conjunction with others to go astray in belief, \textit{lead away with}
							\item[2] to adjust to a condition or circumstance, \textit{accommodate} (\textit{BDAG lists Romans 12:16 under this meaning, but the discussion of meaning \#3 leaves it open that the verse could be referring to people as well---it depends on whether you are taking the \gr{τοῖς ταπεινοῖς} to refer to people or things. A gloss on the verse with meaning \#2 for \gr{συναπάγω} would be ``accommodate yourself to humble ways,'' which I like.})
							\item[3] to join the company of others, \textit{associate with} 
						\end{compactitem}
					\item[LSJ , PDF ] \textbf{}
						\begin{compactitem}
							\item 
						\end{compactitem}
				\end{compactdesc}
				
				\begin{compactdesc}
					\item[Some NT uses] \textbf{}
						\begin{compactdesc}
							\item[]
						\end{compactdesc}
				\end{compactdesc}
		
	% -----
	\subsubsection{Book of Mormon} % *BOM
	% -----
		\label{CONDESCENSION-BOM}
	
		\begin{tabular}{ll}
			Total occurrences in the BOM & 5
		\end{tabular}
		
		\begin{compactdesc}
			\item[1 Nephi 11:16]
			\item[1 Nephi 11:26]
			\item[2 Nephi 4:26]
			\item[2 Nephi 9:53] 
			\item[Jacob 4:7]
		\end{compactdesc}
		
	% -----
	\subsubsection{Doctrine and Covenants} % *DC
	% -----
		\label{CONDESCENSION-DC}
	
		\begin{tabular}{ll}
			Total occurrences in the DC & 0
		\end{tabular}
		
		% \begin{compactdesc}
			% \item[]
		% \end{compactdesc}
		
	% -----
	\subsubsection{Pearl of Great Price} % *PGP
	% -----
		\label{CONDESCENSION-PGP}
	
		\begin{tabular}{ll}
			Total occurrences in the PGP & 1
		\end{tabular}
		
		\begin{compactdesc}
			\item[JS-H Note 5]
		\end{compactdesc}
		
% % ----------------------
% \subsection{Key Phrases} % *KEYPHRASES *PHRASES
% % ----------------------
	% \label{CONDESCENSION-PHRASES}

	% \begin{compactdesc}
		% \item[] \textbf{#times} | list
	% \end{compactdesc}
	
% % ----------------------
% \subsection{Bible Dictionaries} % *DICTIONARIES
% % ----------------------
	% \label{CONDESCENSION-BD}

% % ----------------------
% \subsection{Restoration Usage} % *RESTORATION
% % ----------------------
	% \label{CONDESCENSION-RESTORATION}

	% % -----
	% \subsubsection{Joseph Smith} % *JS
	% % -----
		% \label{CONDESCENSION-JS}
	
		% \begin{compactdesc}
			% \item[{\hyperref[KEY]{Date/Source}}]
		% \end{compactdesc}
	
	% % -----
	% \subsubsection{Brigham Young} % *BY
	% % -----
		% \label{CONDESCENSION-BY}
	
		% \begin{compactdesc}
			% \item[{\hyperref[KEY]{Date/Source}}]
		% \end{compactdesc}
	
% ----------------------
\subsection{Commentary and Studies} % *COMMENTARY *STUDIES
% ----------------------
	\label{CONDESCENSION-COMMENTARY}

	% % -----
	% \subsubsection{Key Points}
	% % -----
		% \label{CONDESCENSION-KEYS}
	
		% \begin{compactitem}
			% \item
		% \end{compactitem}
		
	% -----
	\subsubsection{``Condescension'' in 1 Nephi 11}
	% -----
		\label{CONDESCENSION-1nephi}
		
		The word ``condescension'' appears twice in \hyperref[1NEPHI-11]{1 Nephi 11}, in verses 16 and 26. Here they are:
		
		\begin{quote}
			And he saith unto me: Knowest thou the condescension of God?
		\end{quote}
		
		\begin{quote}
			And the angel said unto me again: Look and behold the condescension of God!
		\end{quote}
		
		Pretty brief verses, so we need the context to understand what is going on. Nephi is seeing an elaboration of Lehi's dream about the tree. He learns about Christ being born with a mortal mother and then in verse 26 he is instructed to ``Look and behold the condescension of God!'' He looks and sees Christ doing a number of different things in the next three verses, including being baptized (27), ministering to people and being thrown out from their society (28), and gathering disciples (29). So the ``condescension of God'' has something to do with Christ doing those things, or actions like those.
		
		Nephi speaks elswhere, in \hyperref[2NEPHI-31-5]{2 Nephi 31:5--7}, about how extraordinary it is that Christ was baptized. So there's a connection to make there between what he sees in chapter 11 and what he later teaches about the baptism in 2 Nephi 31. Maybe what he saw in the baptism really stuck with him.
		
		Another thing to note is that the definitions of \hyperref[SUNAPAGO]{\gr{συναπάγω}} are not that helpful in understanding the term in 1 Nephi. ``Condescension'' only appears once in the KJV NT, and the definitions of \gr{συναπάγω} have to do with accommodating yourself to new circumstances or being willing to associate with people that you normally wouldn't. That isn't what is being described in 1 Nephi 11, where it refers more to Christ giving something much more significant up to come down to Earth. It isn't totally clear to me what he abandons; in \hyperref[ALMA-5-50]{Alma 5:50}, for example, we learn that Christ comes ``in his glory'' (compare \hyperref[1NEPHI-11-28]{1 Nephi 11:28}), so he retained at least some of what he had before. This would be a topic to study in the future.
			
	% -----
	\subsubsection{Condescension and Adam-God}
	% -----
		\label{CONDESCENSION-adamgod}
		
		The word ``condescension'' as it appears in \hyperref[1NEPHI-11]{1 Nephi 11} could be a lens through which you could see Adam-God material there. The ``condescension'' is the losing of the glory, or part of the glory, that accompanies an exalted being. 